% Figure: bifurcation diagram for the bead on a rotating hoop. Below the
% critical speed the bottom is the only stable equilibrium; above it the
% bottom destabilises and two symmetric off-axis equilibria appear.
\begin{figure}[htbp]
\centering
\begin{tikzpicture}[
  axis/.style={draw=engblack, -{Stealth}, thick},
  stable/.style={draw=engblue, very thick},
  unstable/.style={draw=engred, very thick, dashed},
  guide/.style={draw=enggray, dashed, thin},
  label/.style={font=\small},
]
% x = rotation parameter R Omega^2 / g from 0 to 3 (plotted 0..9)
% y = equilibrium angle theta_eq from -pi/2 to pi/2 (plotted -3..3)
\draw[axis] (0,0) -- (9.6,0) node[anchor=west, label] {$R\Omega^2/g$};
\draw[axis] (0,-3.4) -- (0,3.4) node[anchor=north west, label] {$\theta_\text{eq}$};

% critical value at R Omega^2 / g = 1 -> x = 3
\draw[guide] (3,-3.2) -- (3,3.2);
\node[font=\scriptsize, enggray, anchor=south] at (3,3.0) {threshold $R\Omega^2=g$};

% x ticks
\foreach \x/\m in {3/$1$, 6/$2$, 9/$3$} {
  \draw (\x,-0.08) -- (\x,0.08);
  \node[font=\scriptsize, below] at (\x,-0.12) {\m};
}
% y ticks
\node[font=\scriptsize, left] at (-0.12, 0) {$0$};
\node[font=\scriptsize, left] at (-0.12, 3) {$\pi/2$};
\node[font=\scriptsize, left] at (-0.12, -3) {$-\pi/2$};

% Branch theta = 0: stable for x < 3, unstable for x > 3
\draw[stable] (0,0) -- (3,0);
\draw[unstable] (3,0) -- (9,0);
\node[engblue, font=\scriptsize, anchor=south] at (1.5, 0.05) {bottom stable};
\node[engred, font=\scriptsize, anchor=north] at (6, -0.05) {bottom now unstable};

% Off-axis branches: cos(theta_c) = g/(R Omega^2) = 1/p, p = x/3
% theta_c = acos(3/x) for x > 3, in radians; plotted y = theta_c * (3 / (pi/2)) = theta_c * 1.9099
% sample at x = 3..9
\draw[stable] plot[domain=3.05:9, samples=60]
  ({\x}, {1.9099 * acos(3/\x) * 3.14159/180});
\draw[stable] plot[domain=3.05:9, samples=60]
  ({\x}, {-1.9099 * acos(3/\x) * 3.14159/180});
\node[engblue, font=\scriptsize, anchor=south west] at (6.5, 2.0) {$+\theta_c$ stable};
\node[engblue, font=\scriptsize, anchor=north west] at (6.5, -2.0) {$-\theta_c$ stable};

\end{tikzpicture}
\caption[Pitchfork bifurcation of the bead on a rotating hoop]{The
equilibrium angle of a bead on a frictionless vertical hoop rotating at
angular speed $\Omega$, as a function of the dimensionless group
$R\Omega^2/g$. Below the threshold $R\Omega^2=g$ the bottom of the hoop
($\theta=0$) is the only stable equilibrium (solid blue). At the
threshold the bottom loses stability (red dashed) and two symmetric
off-axis equilibria $\pm\theta_c$ with $\cos\theta_c=g/(R\Omega^2)$
branch off and are stable. The shape is the pitchfork bifurcation, the
prototype of symmetry-breaking transitions; the energy method delivers
it from a single line, the stationarity of an effective potential.}
\label{fig:vol03ch05:hoop-bifurcation}
\end{figure}
